% !TEX root = ../../main.tex
\section{Thiết kế các đối tượng nâng cao (Views, Functions \& Triggers)}
\label{sec:3_4_doi_tuong_nang_cao}

Để hiện thực hóa các quy tắc nghiệp vụ đặc thù trong quản lý dòng họ, tự động hóa bảo vệ tính toàn vẹn dữ liệu và tối ưu hóa việc phân tích thống kê đa chiều, hệ thống \textbf{FamilyConnect} tận dụng ngôn ngữ thủ tục mạnh mẽ \textbf{PL/pgSQL} của PostgreSQL để xây dựng các đối tượng CSDL nâng cao.

\subsection{Trigger kiểm soát toàn vẹn Cây phả hệ (Anti-Loop Trigger)}

\subsubsection{Bài toán vòng lặp tổ tiên trong CSDL phân cấp}
Trong mô hình cây gia phả tự tham chiếu qua hai cột \texttt{MaCha} và \texttt{MaMe}, lỗi nghiệp vụ nghiêm trọng nhất là hiện tượng xuất hiện chu trình lặp huyết thống (ví dụ: Thành viên $A$ là cha của $B$, $B$ là cha của $C$, nhưng do thao tác nhầm lẫn, $C$ lại được cập nhật làm cha của $A$). Lỗi này sẽ khiến các câu truy vấn đệ quy \texttt{WITH RECURSIVE} bị rơi vào vòng lặp vô tận và treo hệ thống.

\subsubsection{Cài đặt Hàm và Trigger chặn vòng lặp bằng PL/pgSQL}
Hệ thống cài đặt hàm kiểm tra đệ quy kích hoạt trước mỗi thao tác \texttt{INSERT} hoặc \texttt{UPDATE} trên bảng \texttt{THANH\_VIEN}:

\begin{lstlisting}[language=SQL, caption={Hàm và Trigger chống vòng lặp phả hệ (Anti-Loop Trigger)}, label={lst:trigger_anti_loop}]
-- 1. Ham kiem tra chu trinh to tien - hau due
CREATE OR REPLACE FUNCTION fn_check_genealogy_loop()
RETURNS TRIGGER AS $$
DECLARE
    v_conflict_id uuid;
BEGIN
    -- Truong hop 1: Thanh vien khong the tu lam cha hoac me cua chinh minh
    IF (NEW."MaCha" IS NOT NULL AND NEW."MaCha" = NEW."MaThanhVien") OR
       (NEW."MaMe" IS NOT NULL AND NEW."MaMe" = NEW."MaThanhVien") THEN
        RAISE EXCEPTION 'Loi toan ven: Thanh vien khong the tu lam cha hoac me cua chinh minh (ID: %)', NEW."MaThanhVien";
    END IF;

    -- Truong hop 2: Kiem tra neu MaCha hoac MaMe moi la hau due cua chinh thanh vien nay
    IF (TG_OP = 'UPDATE') THEN
        IF (NEW."MaCha" IS DISTINCT FROM OLD."MaCha" AND NEW."MaCha" IS NOT NULL) OR
           (NEW."MaMe" IS DISTINCT FROM OLD."MaMe" AND NEW."MaMe" IS NOT NULL) THEN
            
            -- Truy van de quy tim tat ca cac doi hau due cua thanh vien hien tai
            WITH RECURSIVE HauDueCTE AS (
                SELECT "MaThanhVien"
                FROM "THANH_VIEN"
                WHERE "MaCha" = NEW."MaThanhVien" OR "MaMe" = NEW."MaThanhVien"
                
                UNION ALL
                
                SELECT c."MaThanhVien"
                FROM "THANH_VIEN" c
                INNER JOIN HauDueCTE h ON (c."MaCha" = h."MaThanhVien" OR c."MaMe" = h."MaThanhVien")
            )
            SELECT "MaThanhVien" INTO v_conflict_id
            FROM HauDueCTE
            WHERE "MaThanhVien" IN (NEW."MaCha", NEW."MaMe")
            LIMIT 1;

            IF FOUND THEN
                RAISE EXCEPTION 'Phat hien vong lap pha he: Cha/Me chi dinh (ID: %) la hau due cua thanh vien (ID: %)!', 
                    v_conflict_id, NEW."MaThanhVien";
            END IF;
        END IF;
    END IF;

    RETURN NEW;
END;
$$ LANGUAGE plpgsql;

-- 2. Dang ky Trigger thuc thi truoc khi ghi du lieu
CREATE OR REPLACE TRIGGER trg_anti_loop_thanh_vien
BEFORE INSERT OR UPDATE OF "MaCha", "MaMe"
ON "THANH_VIEN"
FOR EACH ROW
EXECUTE FUNCTION fn_check_genealogy_loop();
\end{lstlisting}

\subsection{Hàm duyệt cây phả hệ đệ quy (Recursive CTE Traversal)}

Hàm \texttt{fn\_lay\_cay\_pha\_he\_hau\_due} nhận mã định danh của một thành viên tổ tiên và độ sâu cần kết xuất (\texttt{p\_max\_depth}), trả về danh sách toàn bộ các thế hệ con cháu kèm bậc thế hệ tương đối và chuỗi đường dẫn phân cấp:

\begin{lstlisting}[language=SQL, caption={Hàm duyệt cây con cháu đệ quy (Recursive CTE)}, label={lst:func_recursive_tree}]
CREATE OR REPLACE FUNCTION fn_lay_cay_pha_he_hau_due(
    p_to_tien_id uuid,
    p_max_depth int DEFAULT 10
)
RETURNS TABLE (
    "MaThanhVien" uuid,
    "HoVaTen" varchar,
    "GioiTinh" varchar,
    "TheHe" int,
    "BacSau" int,
    "MaCha" uuid,
    "MaMe" uuid,
    "DuongDanPhaHe" text
) AS $$
BEGIN
    RETURN QUERY
    WITH RECURSIVE CayPhaHeCTE AS (
        -- Neo de quy: Lay ban ghi to tien goc
        SELECT 
            t."MaThanhVien",
            t."HoVaTen",
            t."GioiTinh",
            t."TheHe",
            0 AS "BacSau",
            t."MaCha",
            t."MaMe",
            t."HoVaTen"::text AS "DuongDanPhaHe"
        FROM "THANH_VIEN" t
        WHERE t."MaThanhVien" = p_to_tien_id
        
        UNION ALL
        
        -- Buoc de quy: Tim con cai truc tiep
        SELECT 
            c."MaThanhVien",
            c."HoVaTen",
            c."GioiTinh",
            c."TheHe",
            p."BacSau" + 1,
            c."MaCha",
            c."MaMe",
            p."DuongDanPhaHe" || ' -> ' || c."HoVaTen"
        FROM "THANH_VIEN" c
        INNER JOIN CayPhaHeCTE p ON (c."MaCha" = p."MaThanhVien" OR c."MaMe" = p."MaThanhVien")
        WHERE p."BacSau" < p_max_depth
    )
    SELECT * FROM CayPhaHeCTE ORDER BY "BacSau" ASC, "TheHe" ASC;
END;
$$ LANGUAGE plpgsql;
\end{lstlisting}

\subsection{Khung nhìn tổng hợp Thống kê và Báo cáo (Views)}

\subsubsection{Khung nhìn \texttt{v\_thong\_ke\_dong\_ho} (Báo cáo quy mô dòng tộc)}
View cung cấp bức tranh tổng thể về quy mô số lượng chi phái, tổng nhân khẩu nội--ngoại, số người còn sống, số người đã khuất và các di sản ghi nhận:

\begin{lstlisting}[language=SQL, caption={View thống kê tổng quan dòng họ}, label={lst:view_dong_ho}]
CREATE OR REPLACE VIEW "v_thong_ke_dong_ho" AS
SELECT 
    d."MaDongHo",
    d."TenDongHo",
    d."QueQuan",
    COUNT(DISTINCT c."MaChiToc") AS "SoChiToc",
    COUNT(DISTINCT tv."MaThanhVien") AS "TongSoThanhVien",
    COUNT(DISTINCT CASE WHEN tv."TinhTrang" = 'Con song' THEN tv."MaThanhVien" END) AS "SoThanhVienConSong",
    COUNT(DISTINCT CASE WHEN tv."TinhTrang" = 'Da mat' THEN tv."MaThanhVien" END) AS "SoThanhVienDaMat",
    COUNT(DISTINCT ds."MaDiSan") AS "TongSoDiSanLuuTru"
FROM "DONG_HO" d
LEFT JOIN "CHI_TOC" c ON d."MaDongHo" = c."MaDongHo"
LEFT JOIN "THANH_VIEN" tv ON c."MaChiToc" = tv."MaChiToc"
LEFT JOIN "DI_SAN_LICH_SU" ds ON tv."MaThanhVien" = ds."MaThanhVien"
GROUP BY d."MaDongHo", d."TenDongHo", d."QueQuan";
\end{lstlisting}

\subsubsection{Khung nhìn \texttt{v\_danh\_sach\_mung\_tho} (Báo cáo mừng thọ cao niên)}
View tự động tính toán tuổi hiện tại của các cụ cao niên còn sống trong dòng họ dựa trên ngày sinh dương lịch, phân nhóm theo các mốc đại thọ truyền thống (70, 75, 80, 85, 90, 95, 100+ tuổi):

\begin{lstlisting}[language=SQL, caption={View lập danh sách mừng thọ các cụ cao niên hàng năm}, label={lst:view_mung_tho}]
CREATE OR REPLACE VIEW "v_danh_sach_mung_tho" AS
SELECT 
    tv."MaThanhVien",
    tv."HoVaTen",
    tv."GioiTinh",
    tv."NgaySinhDuongLich",
    c."TenChiToc",
    d."TenDongHo",
    DATE_PART('year', AGE(CURRENT_DATE, tv."NgaySinhDuongLich"))::int AS "TuoiHienTai",
    CASE 
        WHEN DATE_PART('year', AGE(CURRENT_DATE, tv."NgaySinhDuongLich")) >= 100 THEN 'Dai Tho (100+)'
        WHEN DATE_PART('year', AGE(CURRENT_DATE, tv."NgaySinhDuongLich")) >= 90 THEN 'Thuong Tho (90-99)'
        WHEN DATE_PART('year', AGE(CURRENT_DATE, tv."NgaySinhDuongLich")) >= 80 THEN 'Thuong Tho (80-89)'
        WHEN DATE_PART('year', AGE(CURRENT_DATE, tv."NgaySinhDuongLich")) >= 70 THEN 'Trung Tho (70-79)'
        ELSE 'Chua den tuoi tho'
    END AS "BacMungTho"
FROM "THANH_VIEN" tv
JOIN "CHI_TOC" c ON tv."MaChiToc" = c."MaChiToc"
JOIN "DONG_HO" d ON c."MaDongHo" = d."MaDongHo"
WHERE tv."TinhTrang" = 'Con song' 
  AND tv."NgaySinhDuongLich" IS NOT NULL
  AND DATE_PART('year', AGE(CURRENT_DATE, tv."NgaySinhDuongLich")) >= 70
ORDER BY "TuoiHienTai" DESC;
\end{lstlisting}
